\section{Experimental protocol}
\label{sec:protocol}

The empirical study is organized as five evidence programs with different
information conditions and different inferential roles.  Phase~A isolates
target-state information on analytic references and records derivative-timing
diagnostics.  A single position trace supplies development-only negative
evidence for unfiltered finite-difference pipelines.  Frozen v3 artifacts
support a synthetic command-profile and runtime audit of the direct method.
A fresh, frozen V4 attempt supplies a same-follower P/PV/PVA comparison but
failed a preregistered validity gate and is non-confirmatory.  Finally,
post-freeze checks cover implementation compatibility and profile semantics,
but are not a new locked comparison.  Results from these programs are not
pooled.

\subsection{Phase~A: analytic information and timing controls}

% CLAIM: C03,C04
All Phase~A runs use one degree of freedom, a control period of
\SI{\ControlPeriodMS}{\milli\second}, and limits
\SI{\VelocityLimit}{\radian\per\second},
\SI{\AccelerationLimit}{\radian\per\second\squared}, and
\SI{\JerkLimit}{\radian\per\second\cubed}.  There is no look-ahead, the
minimum requested duration is one control period, and the interface convention
is \(\mathrm{target}[k]\mapsto\mathrm{output}[k+1]\).  The output of one
\OrdinaryRuckig{} update becomes the current state of the next update.  The
three smooth analytic references---a quadratic with an extremum, a cubic, and
a sinusoid---were pre-specified in the repository before the checked-in
aggregate result.  This provenance does not imply a formal pre-registration:
the generation manifest records a dirty worktree, and per-cycle outputs for most
non-P conditions were not retained.

The controlled target-component comparison uses P,
\([p_k,0,0]\), analytic PV truth, \([p_k,v_k^\star,0]\), and analytic PVA
truth, \([p_k,v_k^\star,a_k^\star]\).  The analytic derivatives are oracle
information, not estimates recoverable online from the position stream.  PV
and PVA truth request the same endpoint position and velocity at every tick;
they differ only in requested endpoint acceleration.  Consequently, their
comparison addresses the recorded position endpoint under this protocol.  It
does not test acceleration-estimator quality, internal-profile quality, or
control effort.

% CLAIM: C05,C08
A second Phase~A block distinguishes derivative formula accuracy from
availability.  Backward differences use history only.  A standard centered
estimate labelled at \(t_k\) uses \(p_{k+1}\) and is therefore an offline,
noncausal diagnostic.  A causal centered variant waits one sample, estimates
the state represented at \(t_{k-1}\), and then propagates that estimate to the
latest sample using the stated finite-difference model.  A separate
next-cycle analytic PVA oracle supplies the state belonging to \(t_{k+1}\) to
the update at \(k\).  It consumes one future sample and is used only as a
timing and indexing sanity control.

Position-following metrics are computed over the common interval beginning
after \PhaseACommonWarmupSamples{} warm-up samples and ending before sample
\PhaseAEvaluationStopIndex.  RMSE and maximum
absolute error compare output position with the reference at the recorded
output index.  The reported best lag is the integer-grid shift, searched over
the registered range, that minimizes aligned position RMSE; positive lag
means that output is late.  It is an end-to-end summary and must not be read
as estimator group delay.  Derivative RMSE uses a separate common interior
that removes estimator-specific endpoint rules.

\subsection{Development CSV}
\label{sec:protocol-csv}

% CLAIM: C06,C07,N01
The position CSV contains \CSVRowCount{} rows.  Only its \texttt{value} column is used:
elapsed time, timestamp, and topic fields are ignored, and every row is
assigned the same \SI{\ControlPeriodMS}{\milli\second} interval.  Evaluation
uses samples
\([\PhaseAEvaluationStartIndex,\CSVRowCount)\).  The trace has no registered velocity or
acceleration truth and is a development input, not an independent locked test
or a robot experiment.

The CSV matrix comprises P and unfiltered finite-difference-derived PV and PVA
pipelines using backward, offline-centered, and delayed causal-centered
derivatives.  Each row retains its causal/future-sample label.  Raw targets are
recorded before a fixed admissibility-projection function; projected targets
are then followed by \OrdinaryRuckig.  Projection is therefore part of the
tested pipeline when it triggers.  Because projection is not crossed as an
independent factor, this design cannot identify separate causal effects for
differentiation and projection.  Raw differentiated acceleration and
projection flags are target diagnostics, not measured plant acceleration or
executed-motion violations.

\subsection{Frozen synthetic v3 evidence}
\label{sec:protocol-v3}

% CLAIM: C10,C13
The v3 result layer is immutable and was not rerun for this paper draft.  It
retains whole-trajectory identities, explicit failure records, checksums, and
independently recomputed bounded summaries.  The locked-test denominator is
\VThreeTrajectoryCount{} complete synthetic trajectories.  The direct condition carrying the
frozen label \texttt{one\_step\_governed\_pva\_direct} contributes
\VThreeDirectCycleCount{} command cycles to the continuous V/A/internal-J,
point-admissibility, reachability, projection, and fallback audit.  Adjacent
sequence consistency has a smaller denominator because each trajectory
contributes one fewer transition than command states.

The registered continuous-constraint test accepts a margin no smaller than
\(\VThreeAuditTolerance\) in the corresponding physical unit.  Thus a recorded margin of
zero passes but is not strictly positive.  The audit concerns exact
constant-jerk command profiles in the synthetic triple-integrator protocol; it
does not include robot dynamics, torque, communication, collision, or
hardware execution.

Runtime uses a separately defined population.  After
\VThreeRuntimeWarmupSamples{} warm-up samples per
trajectory, \VThreeRuntimeCycleCount{} locked-test cycles remain for the same
frozen method label.  End-to-end runtime is compared with the
\SI{\ControlPeriodMS}{\milli\second} control period and summarized by empirical
quantiles, maximum, and deadline-miss count.  These are measurements from one
recorded Python environment, not worst-case execution-time bounds.  The
safety denominator and runtime denominator are never substituted for one
another.

The artifact program records \VThreeRawBundleCount{} independently verified
raw bundles and \VThreeBoundedArtifactCount{} bounded artifacts.  Its strict
acceptance summary contains \VThreeRequiredCriterionCount{} criteria, of which
\VThreeRequiredCriterionPassCount{} passed and
\VThreeRequiredCriterionFailureCount{} failed or were unavailable.  Artifact integrity
supports provenance, not the scientific validity of every comparison.  V3
predates V4 and contains no V4 result; the two frozen evidence layers remain
distinct.

\subsection{Fresh V4 same-follower confirmation attempt}
\label{sec:protocol-v4}

% CLAIM: C14,C16,C17,C19,N03
V4 used fresh experiment identities and seeds with verified zero overlap
against V1--V3 identities.  Its locked test contained six reference families
and \VFourTestTrajectoryCount{} synthetic trajectories.  P, PV, and PVA were
three direct-execution conditions whose declared pipelines differed only in
\texttt{target\_mode}; they shared the estimator, predictor, zero prediction
horizon, governor, direct follower, ideal plant, current-state policy, limits,
and input position stream.  The design therefore compares target components
under the same declared upstream and downstream pipeline, not different
estimators, predictors, or followers.

The inferential unit was the whole trajectory.  The primary PVA-versus-P
analysis required the complete paired denominator, used a
\num{\VFourBootstrapResampleCount}-resample paired percentile bootstrap, and
retained every failed, harmful, or adverse unit.  The locked state machine
allowed one formal test consumption only.  After test visibility, an
experimental-stage defect would freeze the result and require a newly designed
V5 rather than a same-test rerun.  The completed raw experiment and statistical
estimate were retained.  A subsequent report-only repair was restricted to
packaging from the already complete, checksum-verified raw bundle; it neither
generated trajectories nor resumed or reran the experiment.

The preregistered validity rules required complete primary execution, direct
method purity, identical shared information, no hidden replacement or
fallback, and the registered command-profile and continuous-constraint
checks.  Statistical materiality alone could not override these gates.
Likewise, the max-error and lag guardrails and the instrumented full-Python
runtime gate were evaluated separately.  The protocol status was frozen after
test visibility as \texttt{\VFourProtocolStatus}.  The statistical result is
\texttt{\VFourStatisticalClassification}.  Separately, the effective result is
\texttt{\VFourEffectiveClassification}; this classification controls the
claim.  Descriptive observed effects may be reported only together with the
gate failure; confirmatory or causal performance wording is prohibited.

The ordinary-\Ruckig{} matrix is contextual only, and incomplete pairs are
reported as unavailable rather than analyzed as complete cases.  The oracle
matrix substitutes analytic synthetic truth at the same target time and is an
offline, noncausal, nondeployable diagnostic.  Neither context enters the
primary inference.  V4 remains a synthetic ideal-plant protocol and supplies
no real-stream, hardware, deployment, or production-safety evidence.  Full
state transitions, locks, hashes, and evidence dispositions are reported in
\cref{app:v4-confirmation}.

\subsection{Post-freeze checks}

% CLAIM: C11,C12,N03
Post-freeze code distinguishes declared method, native generator, actually
executed command algorithm, shield application, fallback controller, and
command-profile kind.  A compatibility regression reproduces the Phase~A
P-only \OrdinaryRuckig{} result under current code, while profile-aware tests
represent exposed native prefixes as ordered piecewise-constant-jerk
segments.  These checks validate current interfaces and guard against method
identity collapse.  They do not regenerate v3 samples, provide a corrected
locked comparison, rerun V4, or turn V4's observed differences into confirmed
PVA-over-P or PVA-over-PV performance.
